\section{Conclusion}\label{sec:conclusion}
In this paper, we have presented \method, a novel model for implementing 3D object trajectory control in image-to-video synthesis.
Taking depth information combined with K-means clustered points as control signal, our approach captures essential 3D attributes without the need for explicit 3D trajectory estimation. Our user-friendly inference pipeline allows users to input 3D trajectories by simply drawing on 2D images and adjusting point depths, making the synthesis process more accessible. Our model also has certain limitations. First, \method is constrained by the segmentation results of SAM and trajectories provided by the user, and it does not understand physical laws to generate movements of objects without provided trajectory controls. Additionally, since \method was not trained using tracking data, it cannot control the internal pose changes of objects. Finally, the current generation results are limited by the base model SVD. For future work, we aim to extend \method by incorporating more advanced video base models capable of capturing deformable objects and intricate dynamics to better handle non-rigid motions.

\noindent\textbf{Acknowledgements:} This work is supported by the National Key R$\&$D Program of China (No. 2022ZD0160900), the Research Grant Council of the Hong Kong Special Administrative Region under grant number 16203122, Ant Group Research Intern Program, Jiangsu Frontier Technology Research and Development Program (No. BF2024076), and the Collaborative Innovation Center of Novel Software Technology and Industrialization.